\chapter{Conclusion and Future Work}
\label{ch:conclusion}

\section{Conclusions}

This report presents an end-to-end Agentic AI Pipeline for Automated Vulnerability Detection
and Triage. The key findings of this project are:

\begin{itemize}
  \item \textbf{High False Positive Rate of SAST:} Static analysis alone (CodeQL on OWASP Benchmark) produces high recall
        but generates large numbers of false positive alerts, motivating the need for ML-based
        post-processing.
  \item \textbf{Model Efficacy:} We trained and compared Random Forest, Gradient Boosting, and SVM models using a vectorized contextual feature set. Gradient Boosting achieved the highest F1-Score (0.8886) and ROC-AUC (0.9106), demonstrating that ensemble tree models effectively capture complex code context dependencies.
  \item \textbf{Threshold Optimization:} We identified that relying on the default 0.5 decision threshold is sub-optimal. By mathematically optimizing the threshold, we maximized accuracy (e.g., 0.414 for RF, 0.562 for GBM).
  \item \textbf{Agentic Layer:} The pipeline successfully converts high-confidence ML alerts into
        developer-actionable findings with structured remediation suggestions generated by an LLM, reducing the manual burden on developers and closing the loop from detection to patching.
\end{itemize}

The combined result supports the thesis that a layered pipeline—combining static
analysis for coverage, ML classification for precision, and agentic AI for
actionability—is significantly more effective than any single component alone.

\section{Future Work}

Future enhancements to the pipeline include:
\begin{enumerate}
  \item \textbf{Integration of Dynamic Scanning:} Incorporating DAST (Dynamic Application Security Testing) tools like OWASP ZAP into the pipeline to validate static findings dynamically.
  \item \textbf{Real-World Deployment:} Evaluating the end-to-end pipeline in a live production environment with real CI/CD constraints, observing how the agentic patching workflow impacts developer velocity.
\end{enumerate}
